\FloatBarrier

\sporkSection{Grammar and Cornerstones}{sec:Grammar}

The \sporkName{} implementation is embedded as a Python DSL,
but for presentation purposes,
the grammar shown here contains additions from compiler analysis,
clearly marked with ``not in user language''.
We review the fundamental elements of imperative programming,
then discuss the cornerstones of \sporkName{} language design that allow for our instruction usage and synchronization checking.

\sporkMainKey{Conventions:}
Initially-lowercase symbols denote metavariables ranging over the corresponding capitalized set
(e.g. $a_n \in A$, with some exceptions by custom like $\sporks \in \sporkStmt$).
$a^*$ denotes a tuple of $A$ elements.
The set name with a distinguishing subscript denotes a subset of that set
(e.g. $A_n \in \mathcal{P}(A)$, where $\mathcal{P}$ denotes the power set).
Concurrency-specific concepts are denoted with \textgpu{sansSerif} font.

\begin{figure}[h!]
\sporkForceStacktrue
\sporkGrammarStuff
\begin{tabular}{rrll}
\toprule
  $\sporkbool : \sporkBool$ & $\Coloneqq$ & $\sporkFalse | \sporkTrue$ & boolean \\
  $\sporkz : \sporkZ$ & & & integer \\
  $\sporkr : \sporkR$ & & & real number \\
  $\sporkvalue : \sporkValue$ & $\Coloneqq$ & $\sporkbool \,|\, \sporkz \,|\, \sporkr$ & value \\
  $\sporkx : \sporkX$ & & & variable name \\
  $\sporkf : \sporkF$ & & & custom instruction \\
  $\sporkfBehavior : \sporkProc$ & & & instruction behavior (procedure, Figure~\ref{fig:ProcGrammar}) \\
\midrule
  % In the program syntax
  $\sporkextern$ & $\Coloneqq$ & \texttt{sin|cos|tan|exp2|log2}... & extern call \\
  $\sporkcontrol$ & $\Coloneqq$ & \texttt{bool|size|index|stride} & control type \\
  $\sporkprim$ & $\Coloneqq$ & \texttt{f16|bf16|f32|f64|i8|i16}... & primitive data type \\
  $\sporkmem$ & $\Coloneqq$ & \texttt{CudaRmem|CudaGmemLinear}... & memory type \\
  $\sporkspecialWindow$ & $\Coloneqq$ & \texttt{Sm90\_tensorMap(...)} & special window type \\
  $\sporkbarMech$ & $\Coloneqq$ & \texttt{CudaMbarrier|CudaClusterSync}... & barrier mechanism type \\
  $\sporkcollUnit : \sporkCollUnit$ & $\Coloneqq$ & \texttt{cuda\_thread}|\texttt{cuda\_warp}... & collective unit \\
  $\sporkloopMode : \sporkLoopMode$ & $\Coloneqq$ & \texttt{Seq}|\texttt{CudaTasks}|\texttt{CudaThreads}($\sporkcollUnit$) & loop mode \\
  $\sporkwarpName$ & & & warp specialization name \\
\midrule
  % Behind the scenes
  $\sporkmut : \sporkMut$ & $\Coloneqq$ & $\sporkread|\sporkmutate$ & read/mutate tag \\
  $\sporksigma : \sporkSigma$ & & $\sporkSigma \equiv \sporkX \times \mathbb{N}^* \rightharpoonup \sporkValue \cup \{\bot\}$ & value environment (partial function) \\
  $\sporktid : \sporkTid$ & & $\sporkTid \equiv \mathbb{N}^* \times \mathbb{N}$ & thread ID (task, local), Section~\ref{sec:CollAnalysis} \\
  $\sporktidMap : \sporkTidMap$ & & $\sporkTidMap \equiv \sporkSigma \to \mathcal{P}({\sporkTid})$ & thread map, outputs \emph{thread collectives} \\
\bottomrule
\end{tabular}
\caption{Basic Grammar and Types. $\bot$ marks allocated but uninitialized data.}
\label{fig:BasicGrammar}
\sporkCodeRule
\end{figure}


\begin{figure}[h!]
\sporkGrammarStuff
\begin{tabular}{rrll}
\toprule
  $\sporkqtl : \sporkQtl$ & $\Coloneqq$ & \texttt{cuda\_in\_order\_ram\_qual}... & qualitative timeline (\sporkRed{todo table link}) \\
  $\sporkwinAnn : \sporkWinAnn$ & $\Coloneqq$ & $(\sporktidMap, \sporkQtlPre, \sporkqtlInit, \sporkQtlAtom, \sporkQtlMask)$ & window annotation (\textbf{not in user language}) \\
  $\sporkc$ & $\Coloneqq$ & $\sporkexpr | \sporkeLo\texttt{:}\sporkeHi$ & window \textbf{c}oordinate (point $\sporkexpr$ or interval $\sporkeLo\texttt{:}\sporkeHi$) \\
  $\sporkcz$ & $\Coloneqq$ & $\sporkz | \sporkz_\mathrm{lo}\texttt{:}\sporkz_\mathrm{hi}$ & $\sporkZ$-constant window coordinate \\
  $\sporkw$ & $\Coloneqq$ & $\sporkx[\sporkc^*]$ & non-annotated window \\
  $\sporkwz$ & $\Coloneqq$ & $\sporkx[\sporkcz^*]$ & non-annotated $\sporkZ$-constant-indexed window \\
  % Only emph the "optionally" the first time to not be annoying.
  $\sporkwH$ & $\Coloneqq$ & $\sporkw | \sporkw \# \sporkwinAnn$ & window, \emph{optionally} annotated \\
  $\sporkwzH$ & $\Coloneqq$ & $\sporkwz | \sporkwz \# \sporkwinAnn$ & $\sporkZ$-constant-indexed window, optionally annotated \\
  % This should be the only form of scalar window needed.
  % Only needed in the expr/stmt grammar for exposition purposes
  % (better than informal ``expects scalar'' comment).
  % For semantics, it's not really that necessary to constrain to scalars
  % as genericizing to $\sporkwH$ is still correct pattern-matching anyway.
  $\sporkwScalarH$ & $\Coloneqq$ & $\sporkx[\sporkexpr^*] \, | \, \sporkx[\sporkexpr^*] \# \sporkwinAnn$ & scalar window (point coordinates only), optionally annotated \\
\bottomrule
\end{tabular}
%% Scalar allocation access note may need to be emphasized in the right place for a real paper
%% but this is good enough for this staging document for now.
\caption{Window Grammar.
Access scalar allocations $\sporkx$ with the degenerate scalar window $\sporkx[]$ (embedded in Python syntax as plain $\sporkx$).
The window annotation $\sporkwinAnn$ is added by compiler analysis, based on analyzing warp specialization, parallel loops, and (if applicable) the instruction parameter that the window is substituting.
}
\label{fig:WindowGrammar}
\sporkCodeRule
\end{figure}


\sporkMainKey{Mutable Arrays:}
created in a \emph{data allocation statement} (Figure~\ref{fig:stmt}),
or already in-scope as a procedure parameter (Figure~\ref{fig:ProcGrammar}).
The compiler implicitly inserts a \texttt{free} statement after last use;
both the allocation and free statements (for the same $\sporkx$)
have the same multidimensional extents and memory type (Section~\ref{sec:DataLayout}).

Elements in arrays (tensors) are accessed via \emph{windows} (Figure~\ref{fig:WindowGrammar}),
which may reference a single scalar, or a bounding box (slice) of elements.
Crucially, the compiler deduces \emph{window annotations} $\sporkwinAnn$,
which control the trace events generated by accesses through the window.
Due to verbosity, the full small-step semantics are deferred to an appendix
(Section~\ref{sec:Semantics}, Figure~\ref{fig:SyncTraceSemantics}).

The original Exo language already performs type checking and bounds checking,
and defines arithmetic operations (including ``externs'' like exp),
so we do not further discuss these topics.

\sporkMainKey{Loop Nests:} Sequential and parallel loops are distinguished only by $\sporkloopMode$
(Figure~\ref{fig:BasicGrammar})
and, in the case of \texttt{cuda\_threads} loops, $\sporkcollUnit$.
Parallel loops may appear only within CUDA device function launch statements (Figure~\ref{fig:stmt}).
The thread assignment deduced from parallel loops may be further influenced by \emph{warp specialization blocks} (Section~\ref{sec:WarpSpecialization}).

\begin{figure}[h!]
\sporkForceStacktrue
\sporkGrammarStuff
\begin{tabular}{rrll}
\toprule
  $\sporkexpr : \sporkExpr$ & $\Coloneqq$ &
        $\sporkwScalarH$ & scalar read \\
  &$|$& $\sporkvalue$ & constant (boolean, integer, or real) \\
  &$|$& $-\sporkexpr$ & unary negation \\
  &$|$& $\sporkexpr_L \oplus \sporkexpr_R$ & bin-op, $\oplus \in \sporkBinOp$ \\
  &$|$& $\sporkextern(\sporkexpr^*)$ & extern call \\
\bottomrule
\end{tabular}
\begin{tabular}{rll}
\toprule
% Note, we don't say \oplus : \sporkBinOp because \oplus may be used as a generic operator
% placeholder (e.g. = or += in Tr-Mutate), not specifically these BinOps.
$\sporkBinOp \Coloneqq$ & $\{$ \texttt{+},\texttt{-},\texttt{*},\texttt{/},\texttt{\%}, & (arithmetic) \\
  & \texttt{==},\texttt{!=},\texttt{<},\texttt{<=},\texttt{>},\texttt{>=}, & (comparison) \\
  & \texttt{and},\texttt{or} $\}$ & (logical) \\
\bottomrule
\end{tabular}
\caption{Expression Grammar; \texttt{and} and \texttt{or} do \emph{not} short-circuit in Exo-GPU}
\label{fig:expr}
\sporkCodeRule
\end{figure}

\begin{figure}[h!]
\sporkGrammarStuff
\begin{tabular}{rrll}
\toprule
  % Control flow
  $\sporks : \sporkStmt$ & $\Coloneqq$ &
    \texttt{$\sporks_L$;$\sporks_R$} & sequencing \\
  &$|$& \texttt{pass} & pass (skip, no-op) \\
  &$|$& \texttt{if $\sporkexpr$:$\sporks$} & guard \\
  &$|$& \texttt{for $\sporkx$ in $\sporkloopMode$($\sporkeLo,\sporkeHi$):$\sporks$} & sequential or parallel loop \sporkbinder \\
  &$|$& \texttt{with CudaWarps($\sporkz$, $\sporkz$,} & warp specialization block \\
  &   & \texttt{~~~~~~~name=$\sporkwarpName$): $\sporks$} & \\
  &$|$& \texttt{with CudaDeviceFunction(} & cuda device function launch \\
  &   & \texttt{~~~~~~~clusterDim=$\sporkz$,} & \\
  &   & \texttt{~~~~~~~warp\_config=$\sporkwarpConfig^*$): $\sporks$} & ($\sporkwarpConfig$: Figure~\ref{fig:CudaWarps})\\
  &$|$& \texttt{assert $\sporkexpr$} & assertion \\
\midrule
  % Data
  &$|$& \texttt{$\sporkwScalarH$ = $\sporkexpr$} & scalar assignment \\
  &$|$& \texttt{$\sporkwScalarH$ += $\sporkexpr$} & scalar reduction \\
  &$|$& \texttt{$\sporkx$ = $\sporkw$ @ $\sporkspecialWindow$} & window statement \sporkbinder, Section~\ref{sec:TmaLimitations} \\
  &$|$& \texttt{$\sporkf$($\sporkwH^*$)$\sporkfBarH$} & instruction call with \emph{optional} trailing barrier expr $\sporkfBarH$ \\
  &$|$& \texttt{$\sporkx$: $\sporkprim$[$\sporkexpr^*$] @ $\sporkmem$} & data allocation \sporkbinder \\
\midrule
  &$|$& \texttt{$\sporkx:\sporkcontrol$ = $\sporkexpr$} & temporary control value \sporkbinder\ (\textbf{not in user language}) \\
  &$|$& \texttt{free($\sporkx[\sporkexpr^*]$, $\sporkmem$)} & free data (same bounds, $\sporkmem$ as alloc) \\
  &$|$& \texttt{free($\sporkx[\sporkexpr^*]$, $\sporkbarMech$)} & free barriers (same bounds, $\sporkbarMech$ as alloc) \\
  &$|$& \texttt{free($\sporkx[\sporkexpr^*]$, \texttt{None})} & free control value \\
  &   & & (all free stmts \textbf{not in user language}) \\
\midrule
  % Barriers
  &$|$& \texttt{$\sporkx$: barrier[$\sporkexpr^*$] @ $\sporkbarMech$} & barrier allocation \sporkbinder \\
  &$|$& \texttt{Fence($\sporkQtlArrive, \sporkQtlAwait$)} $\#\sporktidMap$ & sync stmt: non-split barrier \\
  &$|$& \texttt{Arrive($\sporkQtlArrive$) >> $\sporkwH^*$} $\#\sporktidMap$ & sync stmt: arrive on split barrier \\
  &$|$& \texttt{Await($\sporkwScalarH$, $\sporkQtlAwait$, $\sporkz_\mathrm{lag}$)} $\#\sporktidMap$ & sync stmt: await on split barrier \\
  &   & & (sync stmt $\#\sporktidMap$ \textbf{not in user language}) \\
\bottomrule
\end{tabular}
\begin{tabular}{rrll}
\toprule
  $\sporkfBarH$ & $\Coloneqq$ & & no trailing barrier expression (blank syntax for instr \texttt{$\sporkf$($\sporkwH^*$)}) \\
  &$|$& \texttt{>> $\sporkwH_\mathrm{bar}$} & trailing barrier expression (for instr \texttt{$\sporkf$($\sporkwH^*$) >> $\sporkwH_\mathrm{bar}$}) \\
  $\sporkfBarZ$ & $\Coloneqq$ & & no trailing barrier expression \\
  &$|$& \texttt{>> $\sporkwz_\mathrm{bar}$} & trailing barrier expression, non-annotated $\sporkZ$-constant-indexed \\
\bottomrule
\end{tabular}
\caption{Statement Grammar;
the $\sporkloopMode$ is embedded in the Python DSL grammar slightly differently than shown here.
The free statement and temporary control value statement are added by compiler analysis;
the latter exists to make it possible to pass arithmetic or logical expressions of control type to instructions (which expect windows $\sporkw$).
The $\sporktidMap$ of synchronization statements are added in a similar manner as window annotations $\sporkwinAnn$.
}
\label{fig:stmt}
\sporkCodeRule
\end{figure}

\begin{figure}[h!]
\sporkGrammarStuff
\begin{tabular}{rrll}
\toprule
  $\sporkformalParam: \sporkFormalParam$ & $\Coloneqq$ & \texttt{$\sporkx$:$\sporkprim$[$\sporkexpr^*$] @ $\sporkmem$} &
  formal data parameter \sporkbinder \\
  &$|$& \texttt{$\sporkx$:$\sporkcontrol$} & formal control parameter \sporkbinder \\
\midrule
  $\sporkp : \sporkProc$ & $\Coloneqq$ &
    \texttt{def($\sporkformalParam^*$):$\sporks_\mathrm{body}$}
    & procedure \\
\bottomrule
\end{tabular}
\caption{Procedure Grammar}
\label{fig:ProcGrammar}
\sporkCodeRule
\end{figure}


\sporkMainKey{Accelerator Instructions:}
Instructions $\sporkf:\sporkF$ may take data or control arguments.
The Hopper-architecture TMA (\texttt{cp.async.bulk}) functionality introduces two complications:
the need to construct named tensor maps,
which are wrapped as \emph{window statements} (Figure~\ref{fig:stmt}) that alias an underlying array,
and the possibility of passing mbarriers directly to instructions,
which are wrapped as \emph{trailing barrier expressions} $\sporkfBarH$ (Figure~\ref{fig:stmt}).

\sporkMainKey{Synchronization:}
We provide a \texttt{Fence} statement for non-split barriers,
and \texttt{Arrive} and \texttt{Await} statements for split barriers (Figure~\ref{fig:stmt}).
These produce $\sporkTrFence$, $\sporkTrArrive$, and $\sporkTrAwait$ trace events, respectively
(Section~\ref{sec:TraceChecks}, Figure~\ref{fig:SyncTraceSemantics}).
The \texttt{Arrive} and \texttt{Await} statements take explicitly allocated \texttt{barrier}-type parameters,
which are of a user-selected barrier mechanism ($\sporkbarMech$: mbarrier, commit group, or cluster sync).

\texttt{Fence} and \texttt{Arrive} are parameterized with an arrive qualitative timeline set $\sporkQtlArrive$.
This specifies the qualitative timelines of the recorded $\sporkTrFence$ or $\sporkTrArrive$ event's inbound observer set,
which filters the kinds of \emph{prior} executed instructions that the synchronization is parallel partially ordered with (Section~\ref{sec:TraceIntro}).
$\sporkQtlAwait$, for \texttt{Fence} and \texttt{Await}, serves a similar purpose for the outbound observer set and \emph{future} executed instructions.
Both sets are embedded in the Python DSL as named \emph{synchronization timeline} objects.
The compiler converts the statement into appropriate synchronization instructions depending on the parameters.
As with windows, the thread map ($\sporktidMap$) is deduced by the compiler.

\sporkMainKey{Cornerstones:}
In order to successfully produce traces,
we rely on \textbf{(1)} \emph{collective analysis} (Section~\ref{sec:CollAnalysis}) to analyze the parallel loops and warp specializations and produce $\sporktidMap$ annotations,
and we rely on \textbf{(2)} \emph{instruction specifications} to produce the qualitative timelines for window annotations (Section~\ref{sec:Qtl}).
% This sentence SUCKS and interrupts the flow but not clear where to place it.
Similar to $\sporktidMap$ mapping statement execution to threads, data allocations (Section~\ref{sec:DataLayout}) also require \emph{distributed memory analysis} to deduce how array shards are mapped across threads.
The non-qualitative-timeline information in the instruction specification (such as collective unit $\sporkcollUnit$) are used in \textbf{(3)} the backend \emph{static instruction usage checks} (Section~\ref{sec:InstrSpec}),
and finally, we perform \textbf{(4)} synchronization checking of the compiled proc,
for a selection of concrete input problem sizes,
using optimized \emph{synchronization environment semantics} (Section~\ref{sec:SyncEnv})
that do checking equivalent to the presented trace semantics.

\FloatBarrier
